% !TEX program = pdflatex
% Thesis chapter: the bridge_tunnel fork_wall environment — seven sequential
% experiments. Self-contained, stock packages only. Figures from
% outputs/thesis_forkwall (copied to paper/figures/thesis_forkwall).
\documentclass[11pt]{article}

\usepackage[margin=1in]{geometry}
\usepackage{graphicx}
\usepackage{booktabs}
\usepackage{amsmath}
\usepackage{xcolor}
\usepackage{caption}
\usepackage{float}
\usepackage[section]{placeins}

\renewcommand{\topfraction}{0.92}
\renewcommand{\floatpagefraction}{0.72}
\captionsetup{font=small, labelfont=bf}
\newcommand{\code}[1]{\texttt{#1}}
\newcommand{\figdir}{figures/thesis_forkwall}

% headline numbers (filled from the final runs)
% generated from the final training runs (mean of last 10 logged iterations)
\newcommand{\succAux}{$0.69$}
\newcommand{\termAux}{$1.00$}
\newcommand{\doorAux}{$0.69$}
\newcommand{\succAuxRep}{$0.70$}
\newcommand{\termAuxRep}{$1.00$}
\newcommand{\doorAuxRep}{$0.70$}
\newcommand{\succNoaux}{$0.99$}
\newcommand{\termNoaux}{$1.00$}
\newcommand{\doorNoaux}{$0.99$}


\title{\textbf{The \code{fork\_wall} Environment:}\\[2pt]
\large A Sequential Experimental Study of a Recurrent Agent's\\
Beliefs, Skills and Decisions in Partially Observed Navigation}
\author{}
\date{}

\begin{document}
\maketitle
\vspace{-2.5em}

\begin{abstract}\noindent
This chapter develops, end to end, the experimental testbed used throughout the
thesis: the \code{bridge\_tunnel} \code{fork\_wall} task, a partially observed
navigation problem in which the map's hidden \emph{category} must be inferred
from terrain evidence and read out in a discrete door choice. We proceed in
seven steps: (\S\ref{sec:env}) the environment; (\S\ref{sec:data}) the
procedural map dataset; (\S\ref{sec:train}) training recurrent PPO agents with
and without an auxiliary belief loss, separating task \emph{completion} from
\emph{decision quality}; (\S\ref{sec:rollouts}) behaviour, including
generalisation to out-of-distribution maps and the three crossing regimes
(build / mine / avoid); (\S\ref{sec:hidden}) what the recurrent state encodes,
via linear probes; (\S\ref{sec:belief}) belief formation along trajectories;
and (\S\ref{sec:plan}) whether the agent's state contains its \emph{future}
crossing strategy before the obstacle is reached --- a minimal test of
model-free planning.
\end{abstract}

% =====================================================================
\section{The environment}
\label{sec:env}

\begin{figure}[H]
  \centering
  \includegraphics[width=\linewidth]{\figdir/fig1_env.png}
  \caption{\textbf{The \code{fork\_wall} task.} \emph{Left:} a full map. The
  agent spawns at the west edge (white dot) and must cross a procedurally
  generated terrain field --- water it can bridge (\emph{build}), rock it can
  tunnel through (\emph{mine}), impassable trees, and open grass --- to reach a
  wall at the east edge. The wall has a single passage; behind it lie
  \emph{two} doors. Only the door matching the map's hidden category pays
  reward: bottom for \code{lakes}, top for \code{rocky}, either for
  \code{balanced}. \emph{Right:} the agent's observation --- a $21{\times}21$
  egocentric crop of tile identities plus scalar features (goal compass,
  facing, step count). The category is never observed directly: it must be
  \emph{inferred} from the terrain statistics seen en route.}
  \label{fig:env}
\end{figure}

\paragraph{Why this design.} The task couples three computations that we want
to study separately in one agent: (i) \emph{navigation} with terrain-dependent
skills (bridging water costs a build action per cell; mining rock likewise);
(ii) \emph{belief formation} --- the map category is a latent variable whose
evidence is spatially distributed and noisy; and (iii) a \emph{discrete
decision} (the door) whose only correct predictor is that belief. Because the
doors are physically identical and spatially separated from the evidence, door
accuracy is a pure behavioural readout of the agent's belief.

\paragraph{Mechanics.} Actions: four movements, \code{build} (place a bridge
tile on adjacent water) and \code{mine} (convert adjacent rock to walkable
ground). Reward: a per-step slack penalty, potential-based shaping toward the
correct door, and a terminal bonus for the category-correct door. Episodes are
capped generously (800 steps) so that timeouts indicate genuine failure to
cross rather than tight scheduling.

% =====================================================================
\section{The map dataset}
\label{sec:data}

\begin{figure}[H]
  \centering
  \includegraphics[width=\linewidth]{\figdir/fig2_dataset.png}
  \caption{\textbf{Procedural maps by category.} Rows: \code{balanced},
  \code{lakes}, \code{rocky}; columns: generator seeds. Categories share one
  generator and differ only in the water/rock coverage fractions, so no single
  tile identifies the category --- classification requires integrating terrain
  statistics. Training samples maps procedurally (fresh seed per episode);
  evaluation uses held-out seeds ($\geq$10\,000).}
  \label{fig:dataset}
\end{figure}

The three categories are deliberately \emph{overlapping} distributions:
\code{balanced} maps contain both lakes and rock formations, and individual
local views of a \code{lakes} map can look rocky. The category is thus a
genuinely latent, spatially extended variable --- the property that makes the
belief analyses of \S\ref{sec:hidden}--\ref{sec:plan} meaningful.

% =====================================================================
\section{Training: auxiliary belief loss vs.\ none}
\label{sec:train}

We train recurrent PPO (GRU core, 128 units; identical architecture and
hyperparameters, seed 2, 6M environment steps) in two variants: with an
\textbf{auxiliary belief loss} (a linear head on the GRU state classifying the
map category, coefficient 0.3, gradients flowing into the recurrent core) and
with \textbf{no auxiliary loss}. The auxiliary head never affects the reward;
it only shapes the representation.

\begin{figure}[H]
  \centering
  \includegraphics[width=\linewidth]{\figdir/fig3_training.png}
  \caption{\textbf{Training curves} (aux-belief seeds 2 and 3, no-aux seed 2).
  (A) mean episode return; (B) success = category-correct door; (C) the
  success \emph{split}: solid = fraction of episodes that terminate at
  \emph{some} door (task completion, no timeout); dashed = correct-door rate
  \emph{among} terminated episodes (decision quality); (D) auxiliary
  belief-head accuracy (aux runs only; chance $1/3$).}
  \label{fig:training}
\end{figure}

The split in panel (C) carries the chapter's two substantive training
observations. First, \emph{completion and decision quality are different
skills with different learning curves}: both agents drive the terminated
fraction (solid) to $1.0$ within the first quarter of training --- crossing
the map, with its terrain-dependent skills, is learned quickly --- while the
correct-door rate among terminated episodes (dashed) develops later and more
slowly. All remaining error is decision error, not failure to cross.

Second, and contrary to the motivating intuition, \emph{the auxiliary belief
loss did not help --- it hurt}. The no-aux agent reaches $\approx 0.99$ door
accuracy; the aux agent plateaus at $\approx 0.69$, and a replicate aux run
with a different seed lands at the same value ($0.70$). Inspecting the aux
agents' behaviour reveals what the plateau is: \textbf{both aux agents take
the bottom door on 100\% of episodes, in every category} --- a constant door
policy, which scores exactly the always-bottom baseline of $2/3$
(\code{balanced} and \code{lakes} correct, \code{rocky} wrong). Meanwhile
their auxiliary heads read the category at $0.77$--$0.85$ accuracy throughout
(panel D). The aux agents therefore \emph{represent} the belief but never
couple it to the door decision; the no-aux agent, which can only earn the last
third of reward by conditioning on the category, escapes the constant-door
local optimum and uses its belief. Representation supervision is no guarantee
of behavioural integration --- an early, purely behavioural instance of the
representation-vs-use dissociation that the later chapters establish causally.
Final performance:

\begin{table}[H]
\centering\small
\begin{tabular}{@{}lccc@{}}
\toprule
 & success & terminated & door-correct $\mid$ terminated \\
\midrule
aux belief (seed 2) & \succAux & \termAux & \doorAux \\
aux belief (seed 3) & \succAuxRep & \termAuxRep & \doorAuxRep \\
no aux (seed 2)     & \succNoaux & \termNoaux & \doorNoaux \\
\bottomrule
\end{tabular}
\caption{End-of-training metrics (mean of last 10 logged iterations).}
\label{tab:final}
\end{table}

% =====================================================================
\section{Behaviour: rollouts, generalisation, regimes}
\label{sec:rollouts}

\begin{figure}[H]
  \centering
  \includegraphics[width=\linewidth]{\figdir/fig4a_rollouts_noaux.png}
  \caption{\textbf{Stochastic rollouts on held-out maps} (30 sampled-policy
  episodes per map; paths overlaid). Titles give success (correct door) and
  wrong-door rates.}
  \label{fig:rollouts}
\end{figure}

\begin{figure}[H]
  \centering
  \includegraphics[width=\linewidth]{\figdir/fig4b_ood_noaux.png}
  \caption{\textbf{Out-of-distribution generalisation.} Success / wrong-door /
  timeout on maps drawn outside the training distribution: larger maps
  ($40{\times}96$ vs.\ $32{\times}64$), quadrupled tree density, and a
  single-cell wall passage.}
  \label{fig:ood}
\end{figure}

\begin{figure}[H]
  \centering
  \includegraphics[width=0.8\linewidth]{\figdir/fig4c_regimes_noaux.png}
  \caption{\textbf{The three crossing regimes, from the agent's eyes.}
  Egocentric observations at the step each regime is expressed: \emph{build}
  (facing water, placing a bridge tile), \emph{mine} (facing rock, tunnelling),
  \emph{avoid} (obstacle ahead, agent detours without using a skill).}
  \label{fig:regimes}
\end{figure}

Three points. First, behaviour is category-appropriate: on \code{lakes} maps
the agent bridges, on \code{rocky} maps it tunnels, and on all maps it
opportunistically detours around small obstacles (the \emph{avoid} regime) ---
a substantial fraction of successful episodes use no skill at all. Second,
generalisation is perfect along the size ($40{\times}96$) and passage-width
axes and degrades only under quadrupled tree density ($0.83$--$0.90$).
Third, the OOD errors are informative about \emph{which} competence breaks:
under dense trees the failures are almost entirely \emph{timeouts}, not
wrong doors --- the navigation problem gets harder, but the belief and its
door read-out remain intact.

% =====================================================================
\section{What the recurrent state encodes}
\label{sec:hidden}

\begin{figure}[H]
  \centering
  \includegraphics[width=\linewidth]{\figdir/fig5_hidden_noaux.png}
  \caption{\textbf{Hidden-state analysis (no-aux agent).} (A) PCA of GRU
  states coloured by the \emph{next} action --- action selection is organised
  in a few directions of state space. (B) Mean activation of the 40 most
  action-selective units by next action: \code{build} and \code{mine} recruit
  distinct unit populations. (C) Linear probes (logistic regression, trained
  on half the episodes, tested on held-out episodes) against
  majority/chance baselines.}
  \label{fig:hidden}
\end{figure}

The probe battery makes the division of labour explicit: from the no-aux
agent's GRU state a linear read-out recovers the map category at $0.83$
(chance $0.33$), its position along the corridor at $0.83$ (chance $0.125$),
the terrain ahead at $0.94$, the next action at $0.76$, and the door it will
eventually choose at $0.85$ --- long before the doors are visible. Crucially,
the category is decodable from an agent that was \emph{never supervised on
it}: the belief emerges from reward alone, because the door decision requires
it.

\begin{figure}[H]
  \centering
  \includegraphics[width=\linewidth]{\figdir/fig5_hidden_aux.png}
  \caption{\textbf{Same analysis for the aux-belief agent.} The auxiliary loss
  yields an even sharper category code ($0.92$ vs.\ $0.83$ probe accuracy) ---
  yet this is the agent that takes the bottom door on every episode
  (\S\ref{sec:train}); its ``final door'' probe is degenerate because the
  behaviour is constant. Decodability of the belief and behavioural use of the
  belief fully dissociate.}
  \label{fig:hidden_aux}
\end{figure}

% =====================================================================
\section{Belief along trajectories}
\label{sec:belief}

\begin{figure}[H]
  \centering
  \includegraphics[width=\linewidth]{\figdir/fig6_belieftraj_noaux.png}
  \caption{\textbf{Trajectories and belief formation (no-aux agent).} Top: single
  stochastic rollouts (colour = time). Bottom: the probe-decoded belief
  $P(\text{category}\mid h_t)$ over the same episodes. The belief is not
  written in one step: it accumulates as terrain evidence enters the view, can
  dip when the local view is ambiguous (a grassy stretch on a \code{lakes}
  map), and is held across the evidence-free corridor before the wall.}
  \label{fig:belieftraj}
\end{figure}

% =====================================================================
\section{Does the agent know its plan?}
\label{sec:plan}

Finally we ask a forward-looking question: when the agent approaches an
obstacle it could either cross with a skill or walk around, does its hidden
state already contain the choice? We build \emph{toy maps} --- a flat grass
corridor with a single water or rock band of controlled extent and gap ---
roll the stochastic policy repeatedly, label each successful episode by the
strategy it eventually expressed (\emph{skill} vs.\ \emph{around}), and train
a probe to predict that label from the hidden state at fixed distances
\emph{before} the obstacle.

\begin{figure}[H]
  \centering
  \includegraphics[width=\linewidth]{\figdir/fig7_planning_noaux.png}
  \caption{\textbf{Strategy is decodable before the obstacle.} (A) Held-out
  probe accuracy for the eventual strategy as a function of distance to the
  obstacle (right = closer). (B) Example toy map with \emph{around} (cyan) and
  \emph{bridge} (pink) episodes from identical starts.}
  \label{fig:planning}
\end{figure}

The curve has a sharp, interpretable structure. While the obstacle is beyond
the agent's $21{\times}21$ view (distance $>10$), accuracy sits at chance ---
there is nothing to know, since the toy corridor contains no earlier evidence.
At $d=10$, \emph{exactly when the obstacle enters the field of view}, accuracy
jumps to $\approx 0.86$, and it climbs to $\approx 0.98$ by contact: the
strategy is largely determined on first sight, several seconds of game time
before the crossing is executed, and is then refined rather than reversed.
The agent does not walk to the obstacle and react; it commits on approach.
We are careful about the inferential limit: decodability before the obstacle
shows the state \emph{contains} the plan, not that the readout is causally
used; the causal analyses of the following chapters take up exactly that
question.

% =====================================================================
\section{Summary}
The chapter established the full experimental loop on \code{fork\_wall}: a
task whose latent category must be inferred (\S\ref{sec:env}--\ref{sec:data});
agents whose completion and decision quality develop on different timescales,
with the surprising twist that auxiliary belief supervision \emph{trapped}
two independent runs in a belief-blind constant-door optimum while the
unsupervised agent coupled its belief to the decision
(\S\ref{sec:train}--\ref{sec:rollouts}); a recurrent state that linearly
encodes belief, action, position and the eventual decision --- with belief
decodability and belief \emph{use} fully dissociating between the two agents
(\S\ref{sec:hidden}--\ref{sec:belief}); and evidence that the state contains
the agent's crossing plan from the moment an obstacle enters view, well before
execution (\S\ref{sec:plan}). These are the ingredients --- and the
measurement tools --- for the causal belief-steering experiments of the next
chapters.

\end{document}
